\documentclass{article}

\usepackage[margin=1in]{geometry}
\usepackage{booktabs}
\usepackage{tabularx}
\usepackage{hyperref}
\usepackage{parskip}
\usepackage{longtable}
\usepackage{array}
\usepackage{enumitem}
\usepackage{pdflscape}

% Ragged-right paragraph column for tables
\newcolumntype{L}[1]{>{\raggedright\arraybackslash}p{#1}}

% Bullet point inside a table cell
\newcommand{\bitem}{\par\noindent\hangindent=1em\hangafter=1%
\makebox[1em][l]{\textbullet}}

% ----- Cross-referencing system ---------------------------------------------
% Every component, hazard and requirement has a unique key. Numbers are worked
% out from where the item appears, and every reference uses the key, so
% reordering the component list, the FMEA rows or the requirement list
% updates every number and link automatically.
%
% Components (Section 3): label each item in the component list with
%   \complabel{key}. \compref{key}{text} links text to it.
%
% Hazards (FMEA table): in the Component cell, write \fmeacomp{key}{name}
%   for the main component, followed by the subcomponent, and give each row
%   \hazdef{key}. Hazards are numbered H<c>-<n>: c is the main component's
%   number in Section 3 and n is the row's position within that component.
%   \hazref{key} links to a hazard.
%
% Requirements (Section 6): start each item with \item[\srdef{key}].
%   Requirements are numbered SR-<n> in list order.
%   \srref{key} links to a requirement.
\newcommand{\complabel}[1]{\phantomsection\label{comp:#1}}
\newcommand{\compref}[2]{\hyperref[comp:#1]{#2}}

\newcounter{hazard}
\newcommand{\fmeacurrent}{}
\renewcommand{\thehazard}{\getrefnumber{comp:\fmeacurrent}-\arabic{hazard}}
\renewcommand{\theHhazard}{\fmeacurrent.\arabic{hazard}}
% Hazard numbering restarts only when the main component changes, so
% several subcomponents of the same component share one sequence.
\newcommand{\fmeacomp}[2]{\def\fmeanext{#1}%
\ifx\fmeanext\fmeacurrent\else\setcounter{hazard}{0}\fi
\gdef\fmeacurrent{#1}\compref{#1}{#2}}
\newcommand{\hazdef}[1]{\leavevmode\refstepcounter{hazard}\label{haz:#1}%
H\thehazard}
\newcommand{\hazref}[1]{\hyperref[haz:#1]{H\ref*{haz:#1}}}

\newcounter{sr}
\newcommand{\srdef}[1]{\refstepcounter{sr}\label{sr:#1}\textbf{SR-\thesr:}}
\newcommand{\srref}[1]{\hyperref[sr:#1]{SR-\ref*{sr:#1}}}

\hypersetup{
    colorlinks=true,       % false: boxed links; true: colored links
    linkcolor=red,          % color of internal links (change box color with linkbordercolor)
    citecolor=green,        % color of links to bibliography
    filecolor=magenta,      % color of file links
    urlcolor=cyan           % color of external links
}

\title{Hazard Analysis\\\progname}

\author{\authname}

\date{}

\input{../Comments.text}
\input{../Common.text}

\begin{document}

\maketitle
\thispagestyle{empty}

~\newpage

\pagenumbering{roman}

\begin{table}[hp]
\caption{Revision History} \label{TblRevisionHistory}
\begin{tabularx}{\textwidth}{llX}
\toprule
\textbf{Date} & \textbf{Developer(s)} & \textbf{Change}\\
\midrule
2026/10/13 & Ahsan Muzammil & Completed section 4 (CA), section 5 and 6 (3 each), section 8 (only my AI uses), and Reflection\\
2026/10/14 & Jake Finlay & Completed Section 3: System Boundaries and
Components, added 3 entries to Section 5: Failure Mode and Effect Analysis,
added 5 saftey requirements to Section 6: Safety and Security Requirements,
added my AI disclosure and reflection \\
\bottomrule
\end{tabularx}
\end{table}

~\newpage

\tableofcontents

~\newpage

\pagenumbering{arabic}

\section{Introduction}

This document presents the hazard analysis for \progname{}, a mobile game
designed to prepare children aged 3--7 for paediatric ophthalmology appointments.
It identifies potential system failures, unsafe behaviours, and other
conditions that could harm users or stakeholders,
and derives the safety and security requirements that eliminate or
reduce those hazards.

A hazard is a system condition that, in combination with conditions in its
environment, has the potential to cause a loss \cite{Leveson2011}. A loss
refers to harm or damage to a stakeholder. For \progname{}, potential losses
include physical or psychological harm to children, breaches of user privacy
or data security, the use of inaccurate information, and failure to prepare
children effectively for their eye appointments.

\section{Scope and Purpose of Hazard Analysis}

The purpose of this hazard analysis is to evaluate risks associated with the
development and use of \progname{}, assess their potential consequences, and
identify measures to prevent or reduce harm. The analysis considers both
technical failures and risks arising from the application's interactions
with young users.

The potential losses considered include:

\begin{itemize}
    \item \textbf{Psychological harm:} Distress or confusion caused by
    age-inappropriate content, misleading information, or unclear instructions.

    \item \textbf{Privacy and security breaches:} Unauthorized disclosure
    of personal information or sensitive user data resulting from unnecessary
    data collection or insecure handling.

    \item \textbf{Ineffective appointment preparation:} Reduced preparedness
    caused by inaccurate information, misleading representations of ophthalmological
    procedures, or inadequate communication of relevant information.

    \item \textbf{Loss of functionality:} Inability to access or use the application
    as intended due to software defects, crashes, or other failures.

    \item \textbf{Loss of user trust:} Reduced confidence in the application
    resulting from unpredictable behaviour, unreliable performance, or unmet user
    expectations.
\end{itemize}

The identified hazards and their potential consequences will guide the selection of
mitigation measures and the formulation of corresponding safety and security requirements.

\section{System Boundaries and Components}

\subsection{System Boundary}

The system boundary contains the game and the device it is played on.

The game includes everything in the release build. This includes the code, content,
assets and any third-party software packaged with it, along with the data it saves
on the device.

The device includes the hardware and operating system that run
the game. The team did not make the devices, but they are inside the boundary
because the game must run correctly on the device.

The boundary ends at the device. The game runs offline after installation and
has no accounts or servers, so everything the system does happens on the
device the child is playing on.

\subsection{Components}

The system contains the following components:

\begin{enumerate}
  \item \complabel{Device}\textbf{The Device:} A phone or tablet running iOS or Android including its operating system, touchscreen,
  speakers or headphones, and storage.

  \item \complabel{Game}\textbf{The Game:} An application installed on the device. The Game can be split into multiple subcomponents:
  \begin{itemize}
    \item \complabel{Gameplay}\textbf{Gameplay:} The levels, the interactions that respond to the child's
    touch input, and scoring.
    \item \complabel{EducationalContent}\textbf{Educational Content:} The facts about each eye procedure and
    the story, characters and illustrations that present them.
    \item \complabel{Quizzes}\textbf{Quizzes:} The questions that check what the child has
    learned, and the feedback given for right and wrong answers.
    \item \complabel{ProgressRewards}\textbf{Progress and Rewards:} Unlocked levels, stars, collectibles,
    cosmetic rewards and player profiles.
    \item \complabel{AudioVisuals}\textbf{Audio and Visuals:} Graphics, animations, effects, music,
    sound effects and narration.
    \item \complabel{SaveData}\textbf{Save Data:} The progress, rewards and settings stored on the
    device.
    \item \complabel{ThirdParty}\textbf{Third-party Software and Assets:} The game engine, libraries,
    art, fonts and sounds included in the game.
  \end{itemize}


\end{enumerate}

\section{Critical Assumptions}

The following assumptions establish the expected conditions under
which \progname{} will operate. Potential failures involving these
conditions will still be considered during the hazard analysis.

\begin{enumerate}

    \item Assume that users have access to a compatible smartphone
    or tablet and possess, or can receive assistance with, the
    basic skills required to interact with a mobile application.

    \textit{Justification:}
    The application is intended for use on consumer mobile devices.
    Since the primary users are young children, their familiarity
    with mobile devices and ability to interact independently may
    vary.

    \item Assume that children will primarily use the application
    at home or in a clinical environment before their ophthalmology
    appointments, where the level of caregiver assistance may vary.

    \textit{Justification:}
    The application is intended to support preparation before
    clinical examinations. The environment and availability of
    caregivers may influence how children interact with the
    application and whether they can complete activities
    independently.

    \item Assume that users will operate the application on 
    consumer mobile devices that meet the minimum hardware 
    and operating system requirements specified by the project.

    \textit{Justification:}
    The application is intended to operate on both iOS and Android
    devices with varying hardware capabilities. Device related
    limitations or failures may affect the availability and
    performance of the application.

    \item Assume that internet connectivity may be unavailable
    or unreliable during normal application use.

    \textit{Justification:}
    Families may use the application in environments with
    limited network connectivity, including clinical waiting
    areas. The application should therefore not depend on
    continuous internet access for its essential preparation
    activities.

    \item Assume that the ophthalmology procedures children 
    experience during appointments may vary depending on the 
    patient, clinician, and clinical setting.

    \textit{Justification:}
    The application prepares children for common ophthalmology
    procedures rather than a specific standardized appointment.
    Variations in clinical practices may affect how closely
    the simulated experience reflects the child's actual
    examination. (The application cannot represent every
    possible variation of a clinical examination)

    \item Assume that qualified clinical professionals will be
    available to provide feedback and review the educational
    content developed for the application.

    \textit{Justification:}
    The development team relies on clinical expertise to ensure
    that the educational content appropriately represents
    ophthalmology procedures. The availability and timing of
    clinical feedback may affect the accuracy and readiness
    of the content.

\end{enumerate}

\section{Failure Mode and Effect Analysis}
\label{sec_FMEA}

\subsection{Method}

The hazards are identified with a Failure Mode and Effects Analysis (FMEA).
FMEA works bottom-up from the components of Section~3: for each component, the
team asked how it could fail and what that failure would do to the child. This
suits the game, because most of its hazards come from a single failure, such
as one piece of wrong content or one unsafe effect, rather than from complex
interactions between components. Failures involving the conditions in the
critical assumptions of Section~4 are still analysed.

The game is meant to be finished in full before the child's visit, since
families do not always know which procedures will be done. A failure that
stops the child from finishing the game therefore stops them from learning
about every procedure that comes after it.

Table~\ref{tab:FMEA} is divided into one section per component of
Section~3. Each row records:

\begin{itemize}
  \item \textbf{Component:} the main component of Section~3 and the
  subcomponent that fails.
  \item \textbf{Failure mode:} how the failure would be observed.
  \item \textbf{Effects of failure:} the consequence for the child or their
  family.
  \item \textbf{Causes of failure:} the reasons the failure could occur.
  \item \textbf{Recommended action:} what will remove the hazard or reduce its
  likelihood or effect.
  \item \textbf{SR:} the safety and security requirements of
  Section~\ref{sec_SR} that the actions lead to.
  \item \textbf{Ref:} the hazard identifier, H$c$-$n$, where $c$ is the main
  component's number in Section~3 and $n$ is the hazard's position within that
  component.
\end{itemize}

\subsection{FMEA Table}

Table~\ref{tab:FMEA} is printed in landscape orientation.

% Column widths as fractions of the landscape text width (after the space
% between columns). The fractions add up to 1.
\newlength{\fmeaw}
\newcolumntype{W}[1]{L{#1\fmeaw}}
\newcommand{\fmeahead}{\toprule
\textbf{Component} & \textbf{Failure Mode} & \textbf{Effects of Failure} &
\textbf{Causes of Failure} & \textbf{Recommended Action} & \textbf{SR} &
\textbf{Ref}\\
\midrule}

\begin{landscape}
\small
\setlength{\fmeaw}{\dimexpr\linewidth-14\tabcolsep\relax}
\begin{longtable}{W{0.12}W{0.13}W{0.16}W{0.18}W{0.26}W{0.09}W{0.06}}
\caption{FMEA table} \label{tab:FMEA}\\
\fmeahead
\endfirsthead
\fmeahead
\endhead
\bottomrule
\multicolumn{7}{r}{continue on next page \ldots}\\
\endfoot
\bottomrule
\endlastfoot

% ----- The Device -----------------------------------------------------------
\fmeacomp{Device}{The Device}
& Device becomes excessively hot during use
& The child may become uncomfortable or stop using the
application; in severe cases, prolonged contact with a
hot device could cause minor injury
& \bitem Device has insufficient cooling
  \bitem Application places excessive load on the device
  \bitem Application continues performing unnecessary
  processing while idle
& \bitem Avoid unnecessary continuous processing
  \bitem Stop resource-intensive activities when they
  are no longer required
  \bitem Test the application on supported devices under
  extended use conditions
& \srref{DeviceResourceUse}
& \hazdef{DeviceOverheating}\\
\addlinespace

& Game experiences poor performance, including
slow response times or dropped frames
& Delayed or inconsistent gameplay may make the
application difficult to use, causing frustration or
preventing the child from following preparation
activities effectively
& \bitem Application requires excessive processing
  \bitem Insufficient memory or processing resources
  \bitem Inefficient rendering or animation
  \bitem Performance degrades during extended use
& \bitem Maintain responsive gameplay and smooth
  visual transitions
  \bitem Optimize processing, rendering, and animations
  \bitem Test gameplay performance on supported devices
  \bitem Verify performance during extended gameplay
& \srref{DevicePerformance}
& \hazdef{PoorPerformance}\\
\addlinespace

% ----- The Game: Gameplay ---------------------------------------------------
\fmeacomp{Game}{The Game}:\newline
\compref{Gameplay}{Gameplay}
& Repeated or accidental touch inputs cause the game
to skip important preparation information
& The child may miss explanations about an
ophthalmology procedure, resulting in confusion
or increased anxiety during the actual examination
& \bitem Multiple touch inputs are processed
during a screen transition
  \bitem Repeated tapping advances through
several instructions unintentionally
  \bitem Gameplay does not distinguish between
a deliberate action and repeated inputs
& \bitem Process only one navigation action
per screen transition
  \bitem Temporarily disable repeated inputs
while transitioning between screens
  \bitem Ensure that essential procedure
information cannot be skipped accidentally
& \srref{PreventAccidentalInput}
& \hazdef{AccidentalInputSkipping}\\
\addlinespace

& Gameplay becomes stuck or enters an unrecoverable state
after an unexpected or invalid player action
& The child may be unable to continue the preparation
activity and may become frustrated or confused
& \bitem Invalid or unexpected input is not handled
  \bitem A gameplay state does not have a valid exit
  \bitem An incomplete action leaves the game in an
  unusable state
  \bitem Error conditions are not tested during gameplay
& \bitem Provide a valid recovery path from every
  gameplay state
  \bitem Handle unexpected or invalid inputs without
  preventing continued play
  \bitem Allow the child to safely restart the current
  activity without losing completed preparation
  \bitem Test gameplay recovery and interruption scenarios
& \srref{GameplayRecovery}
& \hazdef{GameplayStuck}\\
\addlinespace

% ----- The Game: Educational Content ----------------------------------------
\fmeacomp{Game}{The Game}:\newline
\compref{EducationalContent}{Educational Content}
& Procedure shown inaccurately (wrong steps, order or equipment)
& Child's expectations do not match the visit, causing surprise, anxiety and
reduced cooperation during the examination
& \bitem No team member has clinical training
  \bitem Procedures differ between patients, clinicians and clinics
  \bitem Content changed after it was reviewed
& \bitem Have a paediatric eye-care clinician approve all procedure content
  before release
  \bitem Tell the child their doctor may do things a little differently
  \bitem Release only the content version that was approved
& \srref{ClinicianApproval}, \srref{ClinicsDiffer},
  \srref{ApprovedContent}
& \hazdef{InaccurateProcedure}\\
\addlinespace

& Unpleasant sensations left out or played down (``it won't hurt'')
& Child is surprised by stinging drops, a bright light or blurry vision, and
loses trust in their caregivers, clinicians and the game
& \bitem Writers try to reassure the child
  \bitem Sensation information is easy to skip during fast-paced play
& \bitem Describe each sensation honestly before the child sees that part of
  the procedure
  \bitem Never promise that something will not hurt or be uncomfortable
  \bitem Do not let the sensation information be skipped
& \srref{HonestSensations}, \srref{ChildLifeReview}
& \hazdef{SensationsHidden}\\
\addlinespace

& Frightening content (needles, crying, looming equipment,
threatening imagery, or scary story moments)
& Increased fear of the appointment; the child may resist or refuse the
examination
& \bitem Realistic depiction chosen for accuracy
  \bitem Danger or drama added to the story to make the game exciting
  \bitem Art and imagery used without review
& \bitem Have a child-life specialist or child development expert review all
  images, language, and story
  \bitem Keep danger and drama out of the parts of the game that show
  procedures
& \srref{ChildLifeReview}
& \hazdef{Frightening}\\
\addlinespace

& The game does not clearly communicate that
its content is educational rather than medical advice
& Caregivers may mistake simulated procedures or
game feedback for clinical guidance and rely on
the application instead of consulting a clinician
& \bitem The application presents realistic
ophthalmology procedures
  \bitem Educational content is presented
without a clear explanation of its limitations
  \bitem Caregivers are not informed of
the application's intended purpose
& \bitem Display an educational-use notice
during initial setup
  \bitem Require caregiver acknowledgment
before accessing the application
  \bitem Make the notice available in the
caregiver procedure guide
& \srref{MedicalDisclaimer}
& \hazdef{MisinterpretedMedicalPurpose}\\
\addlinespace

% ----- The Game: Quizzes ----------------------------------------------------
\fmeacomp{Game}{The Game}:\newline
\compref{Quizzes}{Quizzes}
& TODO
& TODO
& TODO
& TODO
& TODO
& TODO\\
\addlinespace

% ----- The Game: Progress and Rewards ---------------------------------------
\fmeacomp{Game}{The Game}:\newline
\compref{ProgressRewards}{Progress and Rewards}
& Preparation levels become inaccessible until earlier
levels are completed or sufficient points are earned
& The child cannot access preparation for a procedure
that may be required unexpectedly during an appointment,
resulting in unfamiliarity and possible anxiety
& \bitem Journey progression is implemented with
mandatory level unlocking
  \bitem Access to procedure games depends on
  completion status or accumulated points
& \bitem Make all procedure preparation levels
accessible from the beginning
  \bitem Use points and rewards for motivation
rather than restricting access to educational games
  \bitem Verify that every level can be opened
regardless of previous completion
& \srref{AllLevelsAccessible}
& \hazdef{RestrictedLevelAccess}\\

& Incorrect answers or skipped exercises trigger
negative feedback or penalties
& The child experiences frustration or increased
anxiety and may become reluctant to continue
preparing for ophthalmology examinations
& \bitem Game rewards are awarded based on answer
accuracy rather than activity completion
  \bitem Skipped exercises are treated as failures
  \bitem Gameplay requires correct answers
before allowing progression
& \bitem Provide encouraging feedback for
incorrect responses and skipped exercises
  \bitem Allow children to skip exercises
without losing previously earned rewards
  \bitem Award level completion based on reaching
the ending rather than answering correctly
& \srref{NonPunitiveFeedback}
& \hazdef{PunitiveFeedback}\\
\addlinespace

% ----- The Game: Audio and Visuals ------------------------------------------
\fmeacomp{Game}{The Game}:\newline
\compref{AudioVisuals}{Audio and Visuals}
& Audio, animations, or visual effects continue
playing after the child exits a preparation level
& The child may become confused or distressed
by sounds or visuals from an activity they
have already left
& \bitem Audio playback is not stopped
when the level is exited
  \bitem Ongoing animations or visual effects
are not terminated during screen transitions
  \bitem The exit function changes screens
without stopping active game components
& \bitem Stop all level-specific audio,
animations, and visual effects upon exit
  \bitem Ensure that no level-specific
content continues after returning to
the journey map
  \bitem Test exiting during active
audio playback and animations
& \srref{StopAudioVisualsOnExit}
& \hazdef{AudioVisualsContinueOnExit}\\
\addlinespace

& Sudden loud audio, excessive volume, or rapidly
flashing visual effects cause discomfort or distress
& The child may experience sensory discomfort,
headaches, startle responses, or difficulty
continuing to use the application
& \bitem Audio is played at an unnecessarily high volume
  \bitem Sound effects begin suddenly without a gradual transition
  \bitem Visual effects change or flash rapidly
  \bitem Sensory effects are not tested on supported devices
& \bitem Keep audio at a comfortable volume and avoid
unnecessarily loud sound effects
  \bitem Avoid rapidly flashing or excessively repetitive
visual effects
  \bitem Test audio and visual effects during normal gameplay
  \bitem Allow non-essential audio to be muted or reduced
& \srref{SafeAudioVisuals}
& \hazdef{UnsafeAudioVisuals}\\
\addlinespace

% ----- The Game: Save Data --------------------------------------------------
\fmeacomp{Game}{The Game}:\newline
\compref{SaveData}{Save Data}
& Saved progress becomes corrupted or is reset
& The child loses completed levels, rewards, or
settings and may become frustrated or believe
that previously completed preparation has been lost
& \bitem Application terminates while saving
  \bitem Invalid save data is written
  \bitem Device storage becomes unavailable
  \bitem Changes to the save-data format are
not handled correctly
& \bitem Validate save data before loading it
  \bitem Handle interrupted or invalid saves
without crashing
  \bitem Preserve valid progress when a save
operation fails
  \bitem Test save and load behaviour under
interrupted conditions
& \srref{ReliableSaveData}
& \hazdef{SaveDataCorruption}\\
\addlinespace

& Save data contains unnecessary personal information
or is accessible to other applications
& Information about the child could be exposed,
creating a privacy or security risk
& \bitem Personal information is stored despite
not being required for gameplay
  \bitem Save files are stored without appropriate
device protections
  \bitem Debug information is included in release
save data
& \bitem Store only information required for
offline gameplay
  \bitem Do not store names, contact information,
or other unnecessary personal information
  \bitem Protect stored game data using the
platform's application storage mechanisms
& \srref{MinimalSaveData}
& \hazdef{SaveDataPrivacy}\\
\addlinespace

% ----- The Game: Third-party Software and Assets ----------------------------
\fmeacomp{Game}{The Game}:\newline
\compref{ThirdParty}{Third-party Software and Assets}
& A third-party library or engine component
causes the application to crash or behave incorrectly
& The child may lose progress, be unable to complete
a preparation activity, or receive incomplete
information about an ophthalmology procedure
& \bitem Third-party software contains defects
  \bitem Incompatible library versions
  \bitem Unexpected interaction between dependencies
  \bitem Dependency is updated without sufficient testing
& \bitem Test third-party components as part of
the release build
  \bitem Pin and review dependency versions
  \bitem Test critical gameplay and preparation
flows after dependency changes
  \bitem Remove unused third-party dependencies
& \srref{ThirdPartyValidation}
& \hazdef{ThirdPartyFailure}\\
\addlinespace

& Third-party asset contains inappropriate or
misleading content
& The child may be exposed to frightening,
inappropriate, or inaccurate information that
contradicts the intended preparation experience
& \bitem Asset selected without sufficient review
  \bitem Asset is modified or replaced after review
  \bitem External content does not match the
application's educational requirements
& \bitem Review all third-party assets before release
  \bitem Include third-party assets in the same
clinical and child-life review process as
first-party content
  \bitem Release only reviewed versions of assets
& \srref{ThirdPartyAssetReview}
& \hazdef{ThirdPartyAssetContent}\\
\end{longtable}
\end{landscape}

\section{Safety and Security Requirements}
\label{sec_SR}

The following requirements were identified through the Failure Mode and
Effects Analysis of Section~\ref{sec_FMEA}. Each requirement will also be
incorporated into the SRS to maintain traceability between the two
documents.

\begin{itemize}[leftmargin=1.5cm, labelwidth=1.3cm, labelsep=0.2cm,
  align=left, itemsep=1.2em, parsep=0.6em]

\item[\srdef{ClinicianApproval}] \textbf{Clinical Approval of Procedure Content}

Every procedure shown in the game, including its steps, equipment and the
sensations the child may feel, shall be approved by a paediatric eye-care
clinician before the game is released. The approval shall be recorded.

\textit{Rationale:} No team member has clinical training. A procedure shown
wrongly means the child's expectations do not match the clinic, which causes
the anxiety the game aims to prevent.

\textit{Associated hazards:} \hazref{InaccurateProcedure}

\textit{SRS:} TODO

\item[\srdef{ClinicsDiffer}] \textbf{Variation Between Clinics}

The game shall tell the child, in words and pictures suited to their age,
that their doctor may do things a little differently from the game.

\textit{Rationale:} Procedures vary between patients, clinicians and clinics
(Assumption~5). Telling the child this in advance means a difference at the
visit is less likely to surprise or upset them.

\textit{Associated hazards:} \hazref{InaccurateProcedure}

\textit{SRS:} TODO

\item[\srdef{ApprovedContent}] \textbf{Release of Approved Content Only}

The released game shall contain only the version of the educational content
that was approved by the clinical and child-life reviewers. Any change to
approved content shall be reviewed again before release.

\textit{Rationale:} An approval is only meaningful if the content the child
sees is the content that was reviewed. Late edits, such as wording changes or
new art, could reintroduce errors.

\textit{Associated hazards:} \hazref{InaccurateProcedure}

\textit{SRS:} TODO

\item[\srdef{HonestSensations}] \textbf{Honest Description of Sensations}

For every part of a procedure that the child may feel (such as cold, stinging,
a bright light, pressure or blurry vision), the game shall describe the
sensation before showing that part of the procedure. The description shall
not be skippable, and the game shall not state or imply that a procedure will
not hurt or be uncomfortable.

\textit{Rationale:} Clinicians identified stinging eye drops as a major source
of anxiety. A child who has been told about a sensation is not surprised by
it, and does not learn to distrust the people who prepared them.

\textit{Associated hazards:} \hazref{SensationsHidden}

\textit{SRS:} TODO

\item[\srdef{ChildLifeReview}] \textbf{Child-Life Review of Content}

All images, sounds, language and story in the game shall be reviewed by a
child-life specialist or child development expert before release, confirming
that they are suitable for children aged three to seven, contain no
frightening material, and make no promise that a procedure will not be
uncomfortable.

\textit{Rationale:} Child-life experts judge emotional suitability, which the
clinical review does not cover. This matters more in a game, where story and
excitement are added to the medical content.

\textit{Associated hazards:} \hazref{SensationsHidden}, \hazref{Frightening}

\textit{SRS:} TODO


\item[\srdef{AllLevelsAccessible}]
\textbf{All Preparation Levels Must Remain Accessible}

The game shall allow the child to access every
ophthalmology preparation level regardless of
whether previous levels have been completed.
Journey points, completion status, and rewards
shall not restrict access to preparation games.

\textit{Rationale:}
The procedures required during an ophthalmology
appointment may change unexpectedly. Restricting
access to later levels could prevent children from
preparing for an examination they may need.
Keeping all levels accessible allows children
to explore the procedures in any order.

\textit{Associated hazards:}
\hazref{RestrictedLevelAccess}

\textit{SRS:} TODO


\item[\srdef{MedicalDisclaimer}]
\textbf{Medical Disclaimer and Caregiver Acknowledgment}

The game shall inform caregivers that the
application is intended solely as an educational
preparation aid and does not replace professional
medical advice, diagnosis, or treatment.
The game shall require caregiver acknowledgment
of this notice before first use and make the
notice accessible from the caregiver procedure guide.

\textit{Rationale:}
The game introduces children to ophthalmology
procedures through interactive simulations.
Without a clear explanation of its limitations,
caregivers may mistake the educational content
for professional medical guidance. Requiring
acknowledgment communicates the application's
intended purpose before use.

\textit{Associated hazards:}
\hazref{MisinterpretedMedicalPurpose}

\textit{SRS:} TODO


\item[\srdef{NonPunitiveFeedback}]
\textbf{Encouraging Feedback Without Penalties}

The game shall allow children to skip individual
exercises without losing previously earned points,
coins, or cosmetic rewards. Incorrect responses
and skipped exercises shall not trigger negative
feedback or prevent the child from completing
the preparation level.

\textit{Rationale:}
The application aims to reduce anxiety associated
with ophthalmology examinations. Penalizing children
for incorrect answers or skipped exercises may
cause frustration or discourage participation.
Non-punitive feedback allows children to explore
procedures at their own pace without pressure
to perform correctly.

\textit{Associated hazards:}
\hazref{PunitiveFeedback}

\textit{SRS:} TODO

\item[\srdef{StopAudioVisualsOnExit}]
\textbf{Termination of Audio and Visual Effects on Exit}

When the child exits a preparation level, the game
shall immediately stop all audio playback, animations,
and visual effects associated with that level.
No level-specific audio or visual activity shall
continue after the child returns to the journey map.

\textit{Rationale:}
Children may choose to exit a preparation activity
at any time. If audio or visual effects continue
after exiting, the child may become confused or
distressed by content from an activity they are
no longer participating in. Stopping these
components ensures that exiting the level
fully ends the activity.

\textit{Associated hazards:}
\hazref{AudioVisualsContinueOnExit}

\textit{SRS:} TODO

\item[\srdef{PreventAccidentalInput}]
\textbf{Prevention of Unintentional Gameplay Actions}

The game shall prevent repeated touch inputs from
triggering multiple navigation actions during a
single screen transition. Essential procedure
information shall not be skipped as a result
of repeated or accidental touch inputs.

\textit{Rationale:}
Young children may repeatedly tap the screen
while interacting with the game. If multiple
inputs are processed unintentionally, the child
may skip important explanations about an
ophthalmology procedure. Preventing accidental
navigation helps ensure children receive the
necessary preparation information.

\textit{Associated hazards:}
\hazref{AccidentalInputSkipping}

\textit{SRS:} TODO
\end{itemize}




\section{Roadmap}

\wss{Which safety requirements will be implemented as part of the capstone timeline?
Which requirements will be implemented in the future?}

\newpage{}

\section{Appendix --- Generative AI Use Disclosure}

\subsection{AI Tools Used}

\begin{itemize}
  \item ChatGPT (GPT-5.6 Luna and GPT-6)
  \item Claude (Claude Opus 5.5, Medium Effort)
\end{itemize}

\subsection{How and Where Used}

ChatGPT and Claude were used as supporting tools during the development and formatting
of this document. Its primary purpose was to assist with \LaTeX{} formatting,
improve wording and clarity, and brainstorm potential ideas relevant to the
hazard analysis.

The following sections were assisted by AI:

\begin{itemize}
  \item \textbf{FMEA Table:} Claude was used to help with the table formatting
  and the creation of a numbering system. This allowed us to quickly link the
  hazards to the security requirements without worrying about the order of
  which items appear first.
  
    \item \textbf{Critical Assumptions:}
    ChatGPT was used to brainstorm potential assumptions regarding
    the application's operating environment, intended users, and
    external dependencies. These suggestions were reviewed and
    refined to ensure they were relevant to the project's scope.

    \item \textbf{Safety and Security Requirements:}
    ChatGPT was used to ensure that additional safety and security
    requirements that were brainstormed didn't already overlap with existing SRS 
    NFRs. 

    \item \textbf{Document Formatting:}
    ChatGPT was used to assist with converting and formatting
    written content into appropriate \LaTeX{} structures,
    including lists, tables, and section organization.
\end{itemize}

\subsection{Verification of AI-Generated Content}

All AI-assisted content was reviewed before being incorporated into the
document. The team checked that the suggested assumptions, hazards, and safety
requirements were relevant to the project and consistent with the existing
Problem Statement and Goals, Development Plan, and intended application
functionality identified in SRS document. AI-generated suggestions were
evaluated against the project's actual scope and the hazard analysis
methodology presented in the course lectures. Suggestions that were
inappropriate, unnecessary, or inconsistent with the team's decisions were
removed or revised. For \LaTeX{} formatting, the team verified that the
document compiled successfully and that the resulting output was formatted as
intended. The team remains responsible for the accuracy, completeness,
consistency, and appropriateness of the final document, regardless of whether
content was generated or modified with AI assistance.

\newpage{}

\section{Appendix --- Reflection}

% \wss{Not required for CAS 741}

% \input{../Reflection.text}

% \begin{enumerate}
%     \item What went well while writing this deliverable?
%     \item What pain points did you experience during this deliverable, and how
%     did you resolve them?
%     \item Which of your listed risks had your team thought of before this
%     deliverable, and which did you think of while doing this deliverable? For
%     the latter ones (ones you thought of while doing the Hazard Analysis), how
%     did they come about?
%     \item Other than the risk of physical harm (some projects may not have any
%     appreciable risks of this form), list at least 2 other types of risk in
%     software products. Why are they important to consider?
% \end{enumerate}

\subsection*{Jake Finlay}

\subsubsection*{What went well while writing this deliverable?}

What went best was that we split up the work before we started writing. Every
team member had their own assigned issues on GitHub. This made it easy to find
what to work on next. It also made it clear who was responsible for each part
of the document.

Other than the pain points mentioned below, everything else went fine. Nothing
else noteworthy to write here.

\subsubsection*{What pain points did you experience during this deliverable,
and how did you resolve them?}

One pain point I experienced during this deliverable was that some sections
depend on others. Some of my teammates were waiting on sections I was working
on. To keep the team moving, I had to finish those sections first. This meant
putting off other work that I had planned to complete.

It was also hard to write content that my teammates had not already covered. I
had to keep checking their pull requests and the PDFs produced by our CI
pipeline. This was the only way to make sure I was not repeating something from
another section. To solve this, I downloaded everyone's PDFs and kept them open
on my second monitor. I read through them before I started writing. After that,
writing my sections was fairly straightforward.

The last pain point was linking the safety and security requirements to the
hazards in the FMEA table. The order of the requirements and hazards was likely
to change as everyone's work was merged. This would cause them to be
mislabelled, or force a teammate to renumber all of their requirements and
hazards by hand. To solve this, I created a unique identifier system. Each
requirement and hazard is defined with its own identifier. Every reference
links to that identifier instead of a number. If the order changes, the numbers
and links update automatically. This made it much easier to link everything
together and get the traceability I wanted. I plan to use the same system in
the SRS when I start working on it.

\subsection*{Ahsan Muzammil}

\subsubsection*{What went well while writing this deliverable?}
    Writing the Critical Assumptions section went well because
    it did not depend on other team members completing their
    sections first. This allowed me to work independently
    and complete my assigned section without delays.
    I was able to identify and document the assumptions
    relevant to the system without having to wait for
    other parts of the report to be finalized.

\subsubsection*{What pain points did you experience during this deliverable,
and how did you resolve them?}
    One challenge was contributing to Section 5 (FMEA Diagram)
    and Section 6 (Additional Safety and Security Requirements).
    These sections depended on information from earlier sections,
    particularly the System Boundary and Components section.
    Since those sections were still being developed, I could
    not contribute to Sections 5 and 6 as early as I had planned.

    To resolve this, I communicated with the teammates
    responsible for the dependent sections and asked them
    to prioritize completing their work so I could proceed
    with my contributions. In the meantime, I focused on
    completing the Critical Assumptions section. This
    experience highlighted the importance of coordinating
    dependent tasks and communicating early to prevent
    delays in future deliverables.

\subsection*{Group Answer(s)}

\subsubsection*{Which of your listed risks had your team thought of before
this deliverable, and which did you think of while doing this deliverable? For
the latter ones (ones you thought of while doing the Hazard Analysis), how did
they come about?}

TODO: to be completed once all of the team's hazards and risks have been
added to the document.

\subsubsection*{Other than the risk of physical harm (some projects may not
have any appreciable risks of this form), list at least 2 other types of risk
in software products. Why are they important to consider?}

Two other types of risk in software products are privacy risks and financial
risks. Neither involves physical harm, but both can seriously affect the people
who use a product and the organizations that build it.

Privacy risks come from the personal data that software collects. Many
products store names, locations, contact information and payment details. If
this data is exposed through a breach or a design flaw, users can face
identity theft, fraud, stalking or embarrassment. The harm falls on the users,
yet they usually have no control over how their data is stored or protected.
A breach can also lead to legal penalties for the organization and a lasting
loss of trust in the product. These risks are especially important for our
project, because our users are young children. Children cannot judge what
information is safe to share, and their data is protected by stricter laws and
app store rules. This is one reason our game is designed to work offline and
to collect no personal information at all.

Financial risks come from software failing to do its job correctly or failing
to stay available. A bug in a payment system or a financial calculation can
cost users or organizations money directly. Downtime can also bring a business
to a halt, since many organizations rely on software for their daily
operations. Because a single piece of software is often used by thousands of
people at once, one failure can lead to large and widespread costs. Financial
risk can also come from the design of a product itself. For example, games
with in-app purchases can lead children to spend money without understanding
it. Considering these risks early allows a team to design them out, rather
than repairing the damage after the product is released.

\newpage

\bibliographystyle{unsrt}
\bibliography{../../refs/References}

\end{document}